\expandafter\ifx\csname CAPETINCLUDEMODE\endcsname\relax
\documentclass[11pt,a4paper]{article}
% preamble.tex — packages et commandes communes a toutes les fiches CAPET SII IE
\input{glyphtounicode}
\pdfgentounicode=1
\usepackage{cmap}
\usepackage[T1]{fontenc}
\usepackage[utf8]{inputenc}
\usepackage[french]{babel}
\usepackage[a4paper, top=1.6cm, bottom=1.6cm, left=1.4cm, right=1.4cm]{geometry}
\setlength{\headheight}{14.5pt}
\usepackage{amsmath, amssymb, amsfonts}
\usepackage{siunitx}
\sisetup{locale=FR, output-decimal-marker={,}, inter-unit-product=\ensuremath{{}\cdot{}}}
\usepackage{array, booktabs, multirow, tabularx}
\usepackage{xcolor}
\usepackage[most, breakable]{tcolorbox}
\usepackage{multicol}
\usepackage{enumitem}
\usepackage{mathtools}
\usepackage{lmodern}
\usepackage{microtype}
\usepackage{fancyhdr}
\usepackage{tikz}
\usetikzlibrary{positioning,arrows.meta}
\usepackage{lastpage}

% ── Couleurs par domaine ──────────────────────────────────────────────
\definecolor{colElec}{RGB}{26, 82, 148}       % Électrotechnique — bleu foncé
\definecolor{colElecLight}{RGB}{210, 225, 245} % fond bleu clair
\definecolor{colAuto}{RGB}{34, 120, 66}        % Automatique — vert
\definecolor{colMeca}{RGB}{185, 90, 20}        % Mécanique — orange
\definecolor{colTherm}{RGB}{160, 40, 40}       % Thermique — rouge
\definecolor{colQT}{RGB}{232, 244, 233}        % fond question type (vert très léger)
\definecolor{colJury}{RGB}{255, 243, 205}      % fond jury (jaune ambre)
\definecolor{colJuryFrame}{RGB}{200, 130, 0}   % bordure jury
\definecolor{colGrayLight}{RGB}{245, 245, 247} % fond gris neutre
\definecolor{colSeparator}{RGB}{180, 180, 180} % gris séparateur

% ── Style tcolorbox de section principale ────────────────────────────
\tcbset{
  ficheSection/.style={
    breakable,
    enhanced,
    colback=#1!7!white,
    colframe=#1,
    colbacktitle=#1,
    coltitle=white,
    fonttitle=\bfseries\normalsize,
    left=4pt, right=4pt, top=3pt, bottom=4pt,
    before skip=6pt, after skip=6pt,
  }
}

% ── Style encadré question type ──────────────────────────────────────
\tcbset{
  questiontype/.style={
    breakable,
    enhanced,
    colback=colQT,
    colframe=colElec!55!black,
    colbacktitle=colElec!55!black,
    coltitle=white,
    fonttitle=\bfseries\small,
    left=4pt, right=4pt, top=3pt, bottom=3pt,
    before skip=5pt, after skip=5pt,
  }
}

% ── Style encadré jury ───────────────────────────────────────────────
\tcbset{
  jury/.style={
    breakable,
    enhanced,
    colback=colJury,
    colframe=colJuryFrame,
    fonttitle=\bfseries\normalsize,
    left=4pt, right=4pt, top=3pt, bottom=3pt,
  }
}

% ── Mise en page ─────────────────────────────────────────────────────
\pagestyle{fancy}
\fancyhf{}
\renewcommand{\headrulewidth}{0.4pt}
\renewcommand{\footrulewidth}{0.3pt}

\setlist[itemize]{leftmargin=1.2em, itemsep=1pt, topsep=2pt, parsep=0pt}
\setlist[enumerate]{leftmargin=1.2em, itemsep=3pt, topsep=2pt, parsep=0pt}

\setlength{\parskip}{2pt}
\setlength{\columnsep}{10pt}

\def\CAPETBEGINFILE{\begin{document}}
\def\CAPETENDFILE{\end{document}}
\else
\def\CAPETBEGINFILE{}
\def\CAPETENDFILE{}
\fi

\fancyhead[L]{\textbf{\textcolor{colElec}{T21}} - Fonctions de transfert - Diagramme de Bode}
\fancyhead[R]{\small\textit{CAPET SII IE - Session 2026}}
\fancyfoot[C]{\small\thepage\,/\,\pageref{LastPage}}

\CAPETBEGINFILE

% -----------------------------------------------------------------------
%  BLOC EN-TeTE
% -----------------------------------------------------------------------

\begin{tcolorbox}[
  enhanced, breakable,
  colback=colElec,
  colframe=colElec,
  left=6pt, right=6pt, top=5pt, bottom=5pt,
  before skip=0pt, after skip=4pt
]
  {\color{white}\bfseries\large T21 \quad Fonctions de transfert - Diagramme de Bode}\\[2pt]
  {\color{colElec!30!white}\small
    \textbf{Domaine :} electronique / Automatique - Analyse frequentielle \quad
    \textbf{Presence :} \textbf{9/13} sujets (2013-2026) - theme \emph{fondamental}
  }\\[1pt]
  {\color{colElec!20!white}\footnotesize
    Sous-themes : transformee de Laplace - FdT 1\textsuperscript{er}/2\textsuperscript{e} ordre - identification graphique Bode - marges de phase et de gain - correcteur proportionnel - erreur statique
  }
\end{tcolorbox}

% -----------------------------------------------------------------------
%  1. RAPPELS DE COURS
% -----------------------------------------------------------------------

\begin{tcolorbox}[
  ficheSection=colElec,
  title={1. Rappels de cours},
]

% -- 1.1 Transformee de Laplace ----------------------------------------
\textbf{1.1 Transformee de Laplace - rappels essentiels}

\smallskip
Conditions de Heaviside (conditions initiales nulles) :
\[
  \mathcal{L}\!\left\{\frac{dx}{dt}\right\} = p\,X(p)
  \qquad
  \mathcal{L}\!\left\{\int x\,dt\right\} = \frac{X(p)}{p}
  \qquad
  H(p) = \frac{S(p)}{E(p)}
\]

\medskip\hrule\vspace{4pt}

% -- 1.2 FdT 1er ordre ------------------------------------------------
\textbf{1.2 Systeme du 1\textsuperscript{er} ordre - forme canonique}

\[
  H(p) = \frac{K}{1 + \tau\,p}
  \qquad K : \text{gain statique} \qquad \tau : \text{constante de temps}
\]

Reponse indicielle - l'echelon $E_0$ : $s(t) = K\,E_0\,(1 - e^{-t/\tau})$ \quad;\quad $t_{r5\%} \approx 3\tau$

\textbf{Identification sur le Bode :}
\begin{itemize}
  \item Plateau basse frequence $= 20\log K$ dB
  \item Pulsation de coupure = intersection des asymptotes - 0\,dB/dec et $-20$\,dB/dec
\end{itemize}

\medskip\hrule\vspace{4pt}

% -- 1.3 FdT 2e ordre -------------------------------------------------
\textbf{1.3 Systeme du 2\textsuperscript{e} ordre - forme canonique}

\[
  H(p) = \frac{K}{1 + \dfrac{2m}{\omega_0}\,p + \dfrac{p^2}{\omega_0^2}}
\]
\[
  D\% = 100\,e^{-m\pi/\sqrt{1-m^2}} \quad (m<1) \qquad;\qquad t_{r5\%}\cdot\omega_0 = f(m) \text{ (abaque)}
\]

\medskip\hrule\vspace{4pt}

% -- 1.4 Diagramme de Bode asymptotique --------------------------------
\textbf{1.4 Bode asymptotique - regles de trace}

\smallskip
\begin{center}
\renewcommand{\arraystretch}{1.3}
\begin{tabular}{llll}
\toprule
\textbf{Facteur} & \textbf{Gain asymptotique} & \textbf{Phase} & \textbf{Pente} \\
\midrule
Gain pur $K$ & $20\log K$ (constant) & $0^\circ$ & 0 \\
Integrateur $1/p$ & $0\,\mathrm{dB}$ - $\omega=1$ & $-90^\circ$ & $-20\,\mathrm{dB/dec}$ \\
1\textsuperscript{er} ordre $(1+\tau p)^{-1}$ & $0$ puis $-20\,\mathrm{dB/dec}$ - $1/\tau$ & $0^\circ$ puis $-90^\circ$ & \\
2\textsuperscript{e} ordre & $0$ puis $-40\,\mathrm{dB/dec}$ - $\omega_0$ & $0^\circ$ puis $-180^\circ$ & \\
\bottomrule
\end{tabular}
\end{center}
Pour $m < 1/\sqrt{2} \approx 0{,}707$ : pic de resonance $G_{pic} = 20\log\dfrac{1}{2m\sqrt{1-m^2}}$.

\medskip\hrule\vspace{4pt}

% -- 1.5 Marges de stabilite -------------------------------------------
\textbf{1.5 Marges de phase et de gain (FTBO)}

\[
  M_\varphi = 180^\circ + \arg\!\bigl(FTBO(j\omega_u)\bigr) \quad \text{avec } |FTBO(j\omega_u)| = 1 \quad (0\,\mathrm{dB})
\]
\[
  M_G = -20\log|FTBO(j\omega_\pi)| \quad \text{avec } \arg(FTBO(j\omega_\pi)) = -180^\circ
\]

Criteres habituels : $M_\varphi \geq 45^\circ$ (2023, 2026) ; $M_\varphi \geq 90^\circ$ (2024, 2025) ; $M_G \geq 10\,\mathrm{dB}$ (2026).

\textbf{Procedure graphique :}
\begin{enumerate}[leftmargin=1.5em]
  \item Lire $\omega_u$ sur la courbe de gain (croisement 0\,dB).
  \item Lire la phase - $\omega_u$ - $M_\varphi = 180^\circ + \varphi(\omega_u)$.
  \item Lire $\omega_\pi$ sur la courbe de phase (croisement $-180^\circ$).
  \item Lire le gain - $\omega_\pi$ - $M_G = -G_{dB}(\omega_\pi)$.
\end{enumerate}

\medskip\hrule\vspace{4pt}

% -- 1.6 Correcteur proportionnel --------------------------------------
\textbf{1.6 Correcteur proportionnel $C(p) = K_C$}

Ajouter $K_C$ decale le gain de $+20\log K_C$ dB \textbf{sans modifier la phase}.\\
Determination de $K_C$ pour $M_\varphi$ imposee : lire $G_0$ dB - $\omega_u$ souhaite, puis $K_C = 10^{-G_0/20}$.

\end{tcolorbox}

% -----------------------------------------------------------------------
%  2. FORMULAIRE SYNTHeTIQUE
% -----------------------------------------------------------------------

\begin{tcolorbox}[
  ficheSection=colElec,
  title={2. Formulaire synthetique},
]
\small
\begin{multicols}{2}
\renewcommand{\arraystretch}{1.45}
\begin{tabular}{@{}p{4.8cm}l@{}}
\toprule
\textbf{Formule} & \textbf{Usage} \\
\midrule
\multicolumn{2}{@{}l}{\footnotesize\textit{1\textsuperscript{er} ordre}} \\
$H(p) = K/(1+\tau p)$ & Forme canonique \\
$t_{r5\%} \approx 3\tau$ & Rapidite \\
$K = s(\infty)/E_0$ & Identification \\
$\tau$ : lu - $63{,}2\,\%$ de $K\,E_0$ & Identification \\
\midrule
\multicolumn{2}{@{}l}{\footnotesize\textit{2\textsuperscript{e} ordre}} \\
$H(p) = K/(1+2m\,p/\omega_0 + p^2/\omega_0^2)$ & Forme canonique \\
$D\% = 100\,e^{-m\pi/\sqrt{1-m^2}}$ & Depassement \\
\bottomrule
\end{tabular}

\columnbreak

\renewcommand{\arraystretch}{1.45}
\begin{tabular}{@{}p{4.8cm}l@{}}
\toprule
\textbf{Formule} & \textbf{Usage} \\
\midrule
\multicolumn{2}{@{}l}{\footnotesize\textit{Marges de stabilite (FTBO)}} \\
$M_\varphi = 180^\circ + \varphi(\omega_u)$ & Marge de phase \\
$M_G = -G_{dB}(\omega_\pi)$ & Marge de gain \\
$M_\varphi \geq 45^\circ$ ; $M_G \geq 10\,\mathrm{dB}$ & CdC typique \\
\midrule
\multicolumn{2}{@{}l}{\footnotesize\textit{Correcteur et erreur}} \\
$K_C = 10^{-G_0/20}$ & Reglage $K_C$ \\
$\varepsilon_s = E/(1+K_{BO})$ & Erreur statique \\
$\varepsilon_s = 0$ si integrateur BO & Classe 1 \\
\bottomrule
\end{tabular}
\end{multicols}
\normalsize
\end{tcolorbox}

% -----------------------------------------------------------------------
%  3. QUESTIONS TYPES
% -----------------------------------------------------------------------

\begin{tcolorbox}[
  ficheSection=colElec,
  title={3. Questions types - extraites des sujets reels},
]

% -- QT1 -------------------------------------------------------------

% Mini-procedure de resolution type copie concours
\begin{tcolorbox}[questiontype,title={Procedure type copie concours (5 etapes)}]
\begin{enumerate}[leftmargin=1.5em]
  \item \textbf{Qualifier le modele} : hypotheses, domaine de validite, grandeurs d'entree/sortie.
  \item \textbf{Ecrire les lois} : equations de base avec symboles et unites SI.
  \item \textbf{Isoler l'inconnue} : expression litterale, puis application numerique tracee.
  \item \textbf{Controler l'ordre de grandeur} : coherence physique et unite finale.
  \item \textbf{Conclure concours} : phrase courte orientee decision technique.
\end{enumerate}
\end{tcolorbox}

% Visuel de methode : chemin de resolution CAPET
\begin{center}
\begin{tikzpicture}[font=\small, node distance=7mm and 6mm, >=Stealth]
  \node[draw, rounded corners, fill=colElec!8, inner sep=4pt] (n1) {Notions};
  \node[draw, rounded corners, fill=colElec!8, right=of n1, inner sep=4pt] (n2) {Methode};
  \node[draw, rounded corners, fill=colElec!8, right=of n2, inner sep=4pt] (n3) {Application annale};
  \node[draw, rounded corners, fill=colElec!8, right=of n3, inner sep=4pt] (n4) {Verification jury};
  \draw[->, thick, colElec] (n1) -- (n2);
  \draw[->, thick, colElec] (n2) -- (n3);
  \draw[->, thick, colElec] (n3) -- (n4);
\end{tikzpicture}
\end{center}
\begin{tcolorbox}[questiontype,
  title={QT1 - Identification graphique de $H_M(p)$ sur le Bode \hfill\normalfont\footnotesize E2 disc. 2023 (D4)}
]
\textbf{enonce.}
Le diagramme de Bode de $H_M(p)$ montre un plateau BF - $+26{,}4\,\mathrm{dB}$ et deux pulsations de coupure $\omega_a$ et $\omega_b$. Determiner $K_M$, $\omega_a$, $\omega_b$ par lecture graphique. Proposer une expression approchee si $\tau_a \gg \tau_b$.

\tcbline

\textbf{Corrige.}

Gain plateau BF : $20\log K_M = 26{,}4\,\mathrm{dB}$ - $K_M = 10^{26{,}4/20} = 20{,}9$.

$\omega_a$ et $\omega_b$ : intersections des asymptotes (0 dB/dec avec $-20$ dB/dec, puis $-20$ avec $-40$ dB/dec).

Si $\tau_a \gg \tau_b$ - pele dominant : $H_M(p) \approx K_M / (1 + \tau_a p)$.
\end{tcolorbox}

\vspace{4pt}

% -- QT2 -------------------------------------------------------------

\begin{tcolorbox}[questiontype,
  title={QT2 - Marge de phase et reglage $K_C$ \hfill\normalfont\footnotesize E2 disc. 2023 (D7-D8)}
]
\textbf{enonce.}
FTBO tracee pour $K_C = 1$. La phase vaut $-135^\circ$ - $\omega_u$ (croisement 0\,dB). Calculer la marge de phase. Determiner $K_C$ pour obtenir $M_\varphi = 45^\circ$.

\tcbline

\textbf{Corrige.}

Marge de phase courante : $M_\varphi = 180^\circ + (-135e) = 45^\circ$ - condition deje satisfaite pour $K_C = 1$.

Si le gain - la pulsation souhaitee $\omega_u^*$ est $G_0 = +12\,\mathrm{dB}$ (exemple) :
\[
  K_C = 10^{-12/20} = 10^{-0{,}6} \approx 0{,}25
\]
Le decalage du gain de $-20\log K_C = +12\,\mathrm{dB}$ ramene le croisement - $\omega_u^*$ avec $M_\varphi = 45^\circ$.
\end{tcolorbox}

\vspace{4pt}

% -- QT3 -------------------------------------------------------------

\begin{tcolorbox}[questiontype,
  title={QT3 - FTBF du 2\textsuperscript{e} ordre - depassement nul \hfill\normalfont\footnotesize E1 2026 (C7-C8)}
]
\textbf{enonce.}
Montrer que la FTBF d'une boucle d'asservissement se met sous la forme du 2\textsuperscript{e} ordre. Exprimer $\omega_0$ et $m$ en fonction du gain $K_C$. Deduire la valeur de $K_C$ assurant $D\% = 0$.

\tcbline

\textbf{Corrige.}

Identification par identification de la FTBF sous forme canonique - $m$ est une fonction croissante ou decroissante de $K_C$ selon le systeme.

Condition $D\% = 0$ - $m \geq 1$.

Exemple numerique (2026) : $m = 1$ donne $K_C = K_{Clim}$ lu graphiquement sur l'abaque ou calcule analytiquement. Conclure en verifiant aussi $M_G$.
\end{tcolorbox}

\end{tcolorbox}

% -----------------------------------------------------------------------
%  4. ATTENTES DU JURY
% -----------------------------------------------------------------------

\begin{tcolorbox}[
  jury,
  title={4. Attentes du jury - Erreurs recurrentes et points de vigilance},
]

\begin{enumerate}[label=\textbf{\arabic*.}, leftmargin=*, itemsep=5pt]

  \item \textbf{Mettre la FdT sous forme canonique.} Le jury sanctionne l'absence de mise sous forme avec $K$, $\tau$ (ou $m$, $\omega_0$) explicitement identifies.

  \item \textbf{Lire les marges graphiquement.} Ne pas chercher - les calculer analytiquement ; le document reponse fournit le diagramme de Bode - s'en servir.

  \item \textbf{Distinguer FTBO et FTBF.} Les marges se calculent sur la \textbf{boucle ouverte} ; la stabilite s'apprecie sur la \textbf{boucle fermee}.

  \item \textbf{Signe de la marge de phase.} $M_\varphi = 180^\circ + \varphi(\omega_u)$. Si $\varphi = -150^\circ$ alors $M_\varphi = 30^\circ$ (positif - stable).

  \item \textbf{Identification depuis l'essai indiciel.} Lire $\tau$ - $63{,}2\,\%$ (1\textsuperscript{er} ordre). Pour le 2\textsuperscript{e} ordre : lire le depassement et la pseudo-periode pour obtenir $m$ et $\omega_0$.

  \item \textbf{Theme le plus discriminant.} Jury 2024 : - Traite par 79\,\% des candidats, correctement par seulement 31\,\%. -  prevoir un entraenement intensif sur la lecture de Bode.

\end{enumerate}

\medskip
\begin{tcolorbox}[
  colback=colJuryFrame!12!colJury,
  colframe=colJuryFrame,
  left=4pt, right=4pt, top=2pt, bottom=2pt,
  fontupper=\small\itshape
]
Rapport jury 2025 : "Le jury deplore le manque de maetrise des lois basiques de l'electricite (loi des mailles, loi des neuds) necessaires pour etablir les fonctions de transfert. La redaction des demonstrations est exigee."\\
Rapport jury 2023 : "L'automatique est un champ disciplinaire important des SII auquel les candidats ne peuvent se soustraire. Le jury invite - revoir les concepts de base."
\end{tcolorbox}

\end{tcolorbox}

% -----------------------------------------------------------------------
%  PIED - PReSENCE DANS LES SUJETS
% -----------------------------------------------------------------------

\begin{tcolorbox}[
  enhanced, breakable,
  colback=colGrayLight,
  colframe=colSeparator,
  left=4pt, right=4pt, top=3pt, bottom=3pt,
  before skip=4pt, after skip=0pt,
  fonttitle=\small\bfseries,
  title={Presence dans les sujets analyses (2013-2026)},
]
\begin{center}
\footnotesize\renewcommand{\arraystretch}{1.15}
\begin{tabular}{lccccccccccccc}
\toprule
\textbf{Annee} & 13 & 14 & 15 & 16 & 17 & 18 & 19 & 20 & 21 & 22 & 23 & 24 & 25 \\
\midrule
E1 (commune)   & - & - & - & $\checkmark$ & - & - & $\checkmark$ & - & $\checkmark$ & - & $\checkmark$ & - & $\checkmark$ \\
E2 discipl. IE & - & $\checkmark$ & - & - & - & - & - & $\checkmark$ & - & - & $\checkmark$ & $\checkmark$ & $\checkmark$ \\
E2 appl. comm. & - & - & - & - & - & - & - & - & - & - & - & $\checkmark$ & - \\
\bottomrule
\end{tabular}
\end{center}
\smallskip
\footnotesize\textbf{Sous-themes frequents :}
identification graphique FdT (Bode + essai indiciel) - mise sous forme canonique 1\textsuperscript{er}/2\textsuperscript{e} ordre - marges de phase et de gain - correcteur proportionnel $K_C$ - erreur statique (classe 0 vs 1)

\medskip
\textbf{Ordres de grandeur (reperes rapides).}
\begin{itemize}[leftmargin=1.3em]
  \item Reseau BT FR : tension phase-neutre typique \SI{230}{\volt}, tension phase-phase typique \SI{400}{\volt}, frequence \SI{50}{\hertz}.
  \item Courants d'etude frequents en sujet : de quelques \si{\ampere} a quelques dizaines de \si{\ampere}.
  \item En conversion d'energie, un rendement attendu reste generalement entre 0{,}8 et 0{,}98 selon la technologie.
\end{itemize}

\medskip
\textbf{Synthese express (a refaire sans support).}
\begin{enumerate}[leftmargin=1.4em]
  \item Poser les hypotheses du modele et les conventions de signe.
  \item Ecrire 2 a 4 relations cles avec unites.
  \item Realiser un calcul numerique complet et verifier l'ordre de grandeur.
  \item Justifier le resultat par un argument technique de niveau CAPET.
\end{enumerate}

\end{tcolorbox}

\CAPETENDFILE









